\chapter{Methodology}
\label{ch:methodology}

As already mentioned, missing data can arise in countless different ways,
and the choice of methodology depends on the nature of the missingness, the type of data,
and the specific goals of the analysis. We will focus on a subset
of cases, namely those that involve imputation of the missing values with
"information" from the whole dataset (no causal constraint), with in mind
the goal of training a downstream model that is as close as possible to
the one that would have been trained on the complete dataset.
Our choice of data, models, and evaluation methods will be guided
by these constraints. We will nevertheless mention when/how models should
be tweaked to solve adjacent problems.

\section{Data}
\label{sec:data}

As our main baseline, we chose the most studied and reliable dataset in finance:
the end-of-day prices of the S\&P500, in this case fetched with yfinance from
1/1/2010 to 30/5/2026 (potrei cambiare le date). It clearly is not a dataset inflicted
with the curse of missing data; the only structural ones being occasional circuit breaks
and listings/delistings from the index. To obtain a full dataset with no "real"
missing data, we removed each stock that contained at least one missing observation.
Furthermore, we selected a sample of 50 stocks from the index, stratified across industries;
this is a sensible choice because:
\begin{enumerate}
    \item with more 400 stocks the covariance matrix $\Sigma$ quickly becomes unreliable (noisy covariances,
    skewed eigenvalues, unstable inversion\dots) as the number of observations shrinks;
    \item it is easier to check whether the ground data is actually correct;
    \item model training requires considerably less memory, in a time where RAM and GPU prices are 
    skyrocketing;
    \item it is well known that a well constructed portfolio with more than 30 stocks is already
    sufficiently diversified (since volatility scales with the inverse of the square root of the number of securities).
\end{enumerate}

This dataset will serve as the ground truth against which we will compare our results.
In parallel, we will modify it by inserting fake missing observations according to different
rules, as we define in the subsequent sections.
Now an important choice is whether to operate in the prices "space" or the returns space.
There are positives and negatives about each. For one, the raw dataset is prices, and
in practice missing data will come in the form of missing prices.
\section{Models}
\label{sec:models}

\section{Evaluation Methods}
\label{sec:eval}